% ======================================
% Résolution avec les semi-groupes
% ======================================

\subsection{La théorie des semi-groupes}
\begin{frame}
    Soit $X$ un espace de Banach.
    \begin{definition}
        Soit $\big\{T(t)\big\}_{t \geq 0}$ une famille d'opérateurs bornés de $X$ dans $X$, parametrée par un réel positif $t$. On dit que la famille $\big\{T(t)\big\}_{t \geq 0}$ est un \textcolor{astral}{semi-groupe d'opérateurs bornés sur $X$} si 
        \begin{itemize}
            \item[$(i)$] $T(0)=I$,
            \item[$(ii)$] $T(t+s)=T(t)T(s)$ pour tout $t,s \in \bbR_+$ (propriété des semi-groupes). 
        \end{itemize}
    \end{definition}
    \pause
    \begin{definition}
        Soit $\big\{T(t)\big\}_{t \geq 0}$ un semi-groupe d'opérateurs bornés. L'opérateur $A$ défini par 
        \begin{equation*}
            D(A) = \Big\{ x \in X \ \ | \ \ \lim_{h \to 0} \frac{T(h)x-x}{h} \ \text{ existe } \Big\}
        \end{equation*}
        et pour $x \in D(A)$
        \begin{equation*}
            Ax = \lim_{h \to 0} \frac{T(h)x-x}{h}
        \end{equation*}
        est le \textcolor{astral}{générateur infinitésimal du semi-groupe $\big\{T(t)\big\}_{t \geq 0}$}, et $D(A)$ est le domaine de définition de $A$.
    \end{definition}
\end{frame}

\begin{frame}
    \begin{definition}
        Soit $\big\{T(t)\big\}_{t \geq 0}$ un semi-groupe d'opérateurs bornés. On dit que le semi-groupe $\big\{T(t)\big\}_{t \geq 0}$ est \textcolor{astral}{continu au sens fort} si 
        \begin{equation*}
            \forall x \in X, \ \lim_{t \to 0} T(t)x=x
        \end{equation*}
        On dit que le semi-groupe $\big\{T(t)\big\}_{t \geq 0}$ est un \textcolor{astral}{$\scrC^0$ semi-groupe}.
    \end{definition}
    \pause
    \begin{definition}
        Soit $A$ un opérateur de $X$ dans $X$. 
        \begin{itemize}
            \item On appelle \textcolor{astral}{ensemble resolvant de $A$}, l'ensemble $\rho(A)$ des $\lambda \in \bbC$ tels que $\lambda I-A$ est inversible, \textit{i.e.} $(\lambda I-A)^{-1}$ est un opérateur borné de $X$.
            \item La famille $R(\lambda : A):=(\lambda I-A)^{-1}$, où $\lambda \in \rho(A)$, est appelée la \textcolor{astral}{résolvante de $A$}.
        \end{itemize}
    \end{definition}
\end{frame}

\subsection{Cadre théorique et problème de Cauchy}
\begin{frame}
    \begin{itemize}
        \item On va \og regarder \fg{} notre inconnue $[0,a_{\max}] \times \bbR_+ \ni (a,t) \mapsto x(a,t) \in \bbR_+$ comme une fonction $\bbR_+ \ni t \mapsto \tilde{x}(t)$ à valeurs dans un espace vectoriel de fonctions définies sur $\bbR_+$, c'est-à-dire $[0,a_{\max}] \ni a \mapsto \tilde{x}(t)(a) \in \bbR_+$. \pause
        \item On peut réécrire l'équation de McKendrick-Lotka et ainsi définir un opérateur $A$
        \begin{equation*}
            \frac{d}{d t}x(t) = \big(\underbrace{-\mu - \partial_a}_{=:A}\big) x(t)
        \end{equation*} \pause
        \item On obtient alors le problème de Cauchy
        \begin{equation*}
            \left\{\begin{array}{lr}
                \dfrac{d}{d t}x(t) = Ax(t) & \ \ \ t >0 \\
                x(0) = x_0 &
            \end{array}\right.
        \end{equation*} \pause
        \item Il nous faut intégrer la condition au bord du modèle dans le domaine de définition de $A$. Pour ça on pose
        \begin{equation*}
            D(A) := \Big\{ u \in L^1(0,a_{\max}) \ \ | \ \ u'+\mu u \in L^1(0,a_{\max}) \ \text{ et } \ u(0) = \int_0^{+\infty}\beta(a)u(a)\d a \Big\}
        \end{equation*}
    \end{itemize}
\end{frame}

\subsection{Application du théorème de Hille-Yosida}
\begin{frame}
    \begin{proposition}
        $A : D(A) \rightarrow L^1(0,a_{\max})$ est un opérateur fermé.
    \end{proposition}
    \pause
    \begin{proposition}
        $D(A)$ est dense dans $L^1(0,a_{\max})$.
    \end{proposition}
    \pause
    \begin{proposition}
        On a $]\|\beta\|_{L^{\infty}},+\infty[ \ \subset \rho(A)$ et pour chaque $\lambda$ dans cet ensemble la résolvante de $A$ vérifie
        \begin{equation*}
            \vertiii{R(\lambda : A)}_{L^1} \leq \frac{1}{\lambda - \|\beta\|_{L^\infty}}
        \end{equation*}
        où $\vertiii{\cdot}_{L^1}$ désigne la norme d'opérateur de $L^1$.
    \end{proposition}
    \pause
    \begin{itemize}
        \item Les trois propositions précédentes nous assurent que l'opérateur $A$ vérifie les hypothèses du théorème de Hille-Yosida, donc $A$ est le générateur d'un $\scrC^0$ semi-groupe que l'on note $\big\{T(t)\big\}_{t \geq 0}$ vérifiant 
        \begin{equation*}
            \forall t \in \bbR_+, \ \vertiii{T(t)}_{L^1} \leq \exp\big({\|\beta\|_{L^\infty}t}\big)
        \end{equation*}
    \end{itemize}
\end{frame}

\begin{frame}
    \begin{theorem}
        Le problème de Cauchy de McKendrick-Lotka admet une unique solution $x \in \scrC^0\big(\bbR_+, L^1(0,a_{\max})\big)$.
    \end{theorem}
    \pause
    \begin{itemize}
        \item \textbf{Existence et unicité dans $D(A)$.}\pause
        \begin{itemize}
            \item[$\triangleright$] Soit $x_0 \in D(A)$. On pose alors pour tout $t \in \bbR_+$
            \begin{equation*}
                \forall t \in \bbR_+, \ x(t) = T(t)x_0 \in D(A)
            \end{equation*}\pause
            \item[$\triangleright$] Alors avec les propriétés des semi-groupes, on obtient 
            \begin{eqnarray*}
                \frac{d}{d t}x(t) = AT(t)x_0 = Ax(t)
            \end{eqnarray*}\pause
            \item[$\triangleright$] On considère une autre solution $v$ du problème de Cauchy et on pose pour $t \in \bbR_+$ fixé
            \begin{equation*}
                \forall s \in \bbR_+, \ w(s) = T(t-s)v(s)
            \end{equation*}
        \end{itemize}
    \end{itemize}
\end{frame}

\subsection{Identification du semi-groupe}
\begin{frame}
    \begin{itemize}
        \item On a l'expression de la résolvante. Pour tout $\mu \in \ ]\|\beta\|_{L^\infty}, +\infty[$ et 
        \begin{equation*}
            R(\mu : A)x = \int_0^{+\infty}\ex^{-\mu t}T(t)x\d t
        \end{equation*} \pause
        \item En remplaçant par l'expression de la résolvante dans l'équation précédente et en identifiant, on obtient
        \begin{equation*}
            \forall t \in \bbR_+, \ \forall f \in L^1(0,a_{\max}), \ \big(T(t)f\big)(a) := \left\{\begin{array}{cl}\dfrac{\Pi(a)}{\Pi(a-t)}f(a-t) & \text{si } t \leq a, \\
            & \\
            \Pi(a)B(t-a) & \text{si } t > a.
            \end{array}\right.
        \end{equation*} \pause 
        \item La dépendance en $f$ dans l'expression de $T(t)$ dans le cas $t>a$ est implicite.
    \end{itemize}
\end{frame}

\subsection{Comportement asymptotique}
\begin{frame}
    \begin{definition}
        On définit la \textcolor{astral}{borne spectrale de $A$} par 
        \begin{equation*}
            s(A):= \sup\{\Re (\lambda) \in \bbR \ \ | \ \ \lambda \in \sigma(A)\}
        \end{equation*}
        avec $s(A)=-\infty$ si le spectre de $A$ est vide.
    \end{definition}
    \pause
    \begin{definition}
        On dit que le semi-groupe $\big\{T(t)\big\}_{t \geq 0}$ est \textcolor{astral}{irréductible} si $\{0\}$ et $X$ sont les seuls idéaux fermés $\big\{T(t)\big\}_{t \geq 0}$ invariant.
    \end{definition}
    \pause
    \begin{theorem}
        Soit $\big\{T(t)\big\}_{t \geq 0}$ un $\scrC^0$ semi-groupe positif tel que $\omega_{\text{ess}} < \omega_0$. Soit $\varepsilon >0$ donné par le théorème précédent. Si $\big\{T(t)\big\}_{t \geq 0}$ est irréductible et on a $x \in X^+$ et $x^* \in (X')^+$ tels que 
        \begin{equation*}
            \begin{array}{ccccc}
                Ax = s(A)x && A^*x^* = s(A)x^* && \langle x,x^*\rangle =1
            \end{array}
        \end{equation*}
        alors pour tout $\eta \in \ ]0,\varepsilon[$ et pour tout $x_0 \in X$, il existe une constante $M_\eta \geq 1$ telle que 
        \begin{equation*}
            \|\ex^{-s(A)t}T(t)x - \langle x,x_0^*\rangle x_0\| \leq M_\eta \ex^{-\eta t}\|x\|
        \end{equation*}
    \end{theorem}
\end{frame}

\begin{frame} 
    \begin{itemize}
        \item $\omega_{\text{ess}}=-\infty$ et $\omega_0 = s(A) = \alpha^*$ \pause
        \item $x(a) = x(0)\ex^{-\alpha^*a}\Pi(a)$ \pause
        \item $A^*v = -\mu v + v' + \beta v(0)$ \pause
        \item $x^*(a) = x^*(0)\dfrac{\ex^{\alpha^*a}}{\Pi(a)}\Big(1-\displaystyle{\int_0^a}\ex^{-\alpha^*s}K(s)d s\Big)$ \pause 
        \item Pour satisfaire la condition $\langle x,x^*\rangle =1$, on peut par exemple prendre $x(0)=1$ et déterminer $x^*(0)$ satisfaisant l'équation. \pause
    \end{itemize}
    \begin{proposition}
        Le semi-groupe $\big\{T(t)\big\}_{t \geq 0}$ est irréductible si et seulement si 
        \begin{equation*}
            \sup \ ]0,a_{\max}[ \ \setminus \omega = a_{\max}
        \end{equation*}
        où $\omega \subset \ ]0,a_{\max}[$ est l'ouvert constitué des points au voisinage desquels $\beta=0$.
    \end{proposition}
    \pause
    Donc on a par le théorème
    \begin{equation*}
        \lim_{t \to +\infty} \big(T(t)x_0\big)(a) = \lim_{t \to +\infty} \ex^{\alpha^*(t-a)}\Pi(a)\langle x_0,x^*\rangle
    \end{equation*}
\end{frame}